\documentclass{article}
\usepackage[utf8]{inputenc}
\usepackage[T2A]{fontenc}
\usepackage[russian]{babel}
\usepackage{graphicx}
\usepackage{amsmath}
\usepackage{url}
\usepackage{geometry}
\usepackage{setspace}
\usepackage{listings}
\usepackage{xcolor}
\usepackage{longtable}
\usepackage[hidelinks]{hyperref}
\geometry{a4paper, margin=2cm}

\lstset{
  keepspaces=true,
  extendedchars=true,
  basicstyle=\ttfamily\footnotesize,
  keywordstyle=\color{blue!60!black},
  commentstyle=\color{gray},
  stringstyle=\color{green!50!black},
  numbers=left, numberstyle=\tiny\color{gray}, numbersep=6pt,
  breaklines=true, frame=single, language=Python, showstringspaces=false,
  morecomment=[l]{\#}
}

% заглушка для ещё не готовых рисунков
\newcommand{\MaybeImage}[2]{%
  \IfFileExists{#1}{\begin{center}\includegraphics[width=#2\linewidth]{#1}\end{center}}%
  {\begin{center}\fbox{\parbox[c][2.2cm][c]{0.8\linewidth}{\centering\small Рисунок в работе:\\\texttt{\detokenize{#1}}}}\end{center}}%
}
% заглушка для ещё не досчитанного числа в таблице
\newcommand{\Pending}{\textit{--- (в процессе)}}
% заглушка для ненаписанного раздела
\newcommand{\Todo}[1]{\par\medskip\noindent\fbox{\parbox{0.95\linewidth}{\small\textbf{TODO.} #1}}\par\medskip}

\title{Предсказание продуктов протеолиза: честный протокол вложенной кросс-валидации\\ \large для сравнения архитектурных модификаций DeepPeptide}
\author{}
\date{Черновик, \today}

\begin{document}
\maketitle

\begin{abstract}
DeepPeptide~\cite{deeppeptide2023} -- современная модель предсказания продуктов протеолиза (границ зрелых пептидов и пропептидов в белке-предшественнике) на основе эмбеддингов белковых языковых моделей (pLM) и линейно-цепного CRF. Аудит её кода и методологии обнаружил три проблемы: ошибку в подсчёте метрики допуска $\pm3$, смешение выбора архитектуры с итоговой оценкой качества на одних и тех же данных и однородные, не зависящие от позиции в сегменте эмиссии CRF, из-за которых модель не различает начало, середину и конец пептида.

Для решения последней проблемы мы предлагаем две модификации: \emph{boundary-голову}, добавляющую эмиссии для позиционно-фиксированных состояний CRF вблизи границ сегмента, и лёгкий \emph{адаптер} (LayerNorm $\to$ Linear $\to$ GELU), перепроецирующий pLM-эмбеддинг перед основной сетью. Обе оценены протоколом $5\times4$ вложенной кросс-валидации (20 обучений на модель) на пересобранном по методике оригинальной статьи датасете (Swiss-Prot 2026). На эмбеддингах ESM-2 обе модификации дают прирост, складывающийся почти без потерь: $+0.024$ F1 от boundary-головы, $+0.026$ от адаптера, $+0.057$ при совместном использовании (базовая модель $0.556\pm0.028$). Вычисления на эмбеддингах ESM-C 6B на момент написания не завершены; более ранний, менее строгий протокол предсказывает похожую картину, но с заметно большей амплитудой. Более ранние эксперименты без CV-проверки (3Di, bond-лосс, телескопический CRF и другие) вынесены в приложение как наблюдения, которые ещё предстоит подтвердить, а не как доказанные результаты.
\end{abstract}

\section{Введение}

\textbf{Белки} -- это состоящие из аминокислотных остатков биополимеры: они выполняют структурные, каталитические и регуляторные функции в клетках. Неправильно свёрнутые, повреждённые и короткоживущие регуляторные белки, а также чужеродные белки, попавшие в организм, подвергаются \textbf{протеолизу} -- ферментативному расщеплению белковых цепей специализированными протеазами. В результате образуются \textbf{пептиды}: короткие фрагменты, которые нередко обладают собственной биологической активностью и участвуют в сигнальных каскадах и иммунных процессах.

\subsection{Почему эта задача важна}

Протеомные данные обычно содержат полные последовательности белков, а не конечные продукты их расщепления, поэтому предсказание границ образующихся пептидов и пропептидов напрямую влияет на несколько прикладных направлений.

В протеомике такие модели помогают интерпретировать данные масс-спектрометрии: сопоставлять наблюдаемые пептиды с белками-предшественниками, локализовать вероятные места разреза и оценивать, как вариации последовательности меняют паттерн расщепления. При разработке вакцин похожий прогноз нужен уже на этапе проектирования, чтобы понять, какие фрагменты образуются \textit{in vivo} и могут пригодиться при конструировании мультиэпитопных вакцин, а заодно заранее отсеять участки, склонные к нежелательному расщеплению или ускоренной деградации. Третье применение связано с патологией: при нейродегенеративных, онкологических и эндокринных заболеваниях протеолиз часто нарушается, и появляются нетипичные паттерны расщепления. Модель предсказания деградации белков здесь может сопоставлять ожидаемый и наблюдаемый профили в контрольных и патологических образцах и подсказывать гипотезы о том, как накапливаются токсичные или изменённые фрагменты.

\subsection{Недостатки DeepPeptide}
\label{sec:intro-flaws}

DeepPeptide~\cite{deeppeptide2023} -- наиболее близкая по постановке задачи современная модель (раздел~\ref{sec:related}) и отправная точка данной работы. Аудит её кода и методологии обнаружил три проблемы, каждая из которых мотивирует свою часть данной работы.

\textbf{(1) Ошибка в подсчёте метрики.} Функция, реализующая сопоставление предсказанных и истинных сегментов с допуском $\pm3$ (и дающая все P/R/F1 в оригинальной публикации), переиспользует имя переменной цикла: часть верных совпадений в итоге засчитывается не той истинной строке, а полнота и F1 занижаются. Подробности и то, как мы поступили с этой находкой, -- в разделе~\ref{sec:eval-metric}.

\textbf{(2) Смешение выбора архитектуры и итоговой оценки.} В оригинальной методологии внешний фолд вложенной кросс-валидации честно оценивает \emph{конкретную} архитектуру, но сам выбор архитектуры (какие блоки добавлять, какую языковую модель взять) делается по тем же фолдам, на которых затем приводится итоговое качество. Если перебирать варианты и оставлять тот, что лучше показал себя на этих же данных, оценка для лучшего варианта оказывается смещена вверх ровно так, как отбирали. Раздел~\ref{sec:eval} описывает, как мы развели эти два этапа.

\textbf{(3) Одинаковые эмиссии CRF для всех состояний одного типа.} Декодер DeepPeptide -- это CRF с расширенным набором состояний, позволяющим отличить пептид от пропептида и ограничить длину сегмента. Но все состояния внутри одного типа сегмента (скажем, все «пептидные» состояния) получают одну и ту же по-остаточную эмиссию: модель не различает, стоит ли остаток в начале, середине или конце сегмента, хотя сама структура состояний CRF это различие уже кодирует. Раздел~\ref{sec:method} предлагает модификацию, которая эту проблему адресует.

\section{Смежные работы}
\label{sec:related}

Методы решения задачи предсказания протеолитических пептидов делятся на три класса.

\textbf{Модели с привязкой к протеазе.} Простейший класс -- модели, предсказывающие разрез строго для ограниченного набора ферментов, для которых, как правило, известен режущий мотив: \textbf{ProP}~\cite{duckert2004prop} ищет разрезы только семейства PACE/PC у животных и растений, \textbf{PROSPER}~\cite{song2012prosper} и \textbf{PROSPERous}~\cite{song2018prosperous} предсказывают разрезы одного выбранного фермента за раз, развитие последней, \textbf{ProsperousPlus}~\cite{li2023prosperousplus}, устроено так же, \textbf{DeepCleave}~\cite{li2020deepcleave} ограничена каспазами и матриксными металлопротеиназами, а \textbf{Procleave}~\cite{li2020procleave} тоже работает с одним ферментом за раз и вдобавок требует данные о 3D структуре. Ограничение здесь общее для всех: узкая область применимости, тогда как на практике интересен результат совместной работы всех протеаз организма, а не одной.

\textbf{Модели с привязкой к организму или типу пептида.} Отдельный класс моделей специализирован на нейропептидах или иных биологически активных фрагментах конкретного организма. \textbf{NeuroPred}~\cite{southey2006neuropred} и универсальная по видам \textbf{DeepNeuropePred}~\cite{wang2024deepneuropepred} видят разрезы только рядом с несколькими основными аминокислотами; \textbf{NeuroPred-PLM}~\cite{wang2023neuropredplm} вовсе не ищет места разреза, а лишь классифицирует уже готовую последовательность. Рядом стоят модели, ищущие не биоактивные пептиды, а сигнальные подпоследовательности: \textbf{SignalP}~\cite{teufel2022signalp6} и \textbf{TargetP}~\cite{almagro2019targetp} определяют тип сигнального пептида, указывающего, куда белок будет транспортирован.

\textbf{Модели общего назначения на pLM-эмбеддингах (класс DeepPeptide).} \textbf{Белковые языковые модели} (protein language models, pLM) -- это трансформеры, обученные без учителя на больших корпусах последовательностей и выдающие переносимые по-остаточные эмбеддинги: ESM-2~\cite{esm2_2019}, ESM~Cambrian~\cite{esm_cambrian_blog_2024}, Ankh~\cite{ankh2023}. Поверх «замороженных» эмбеддингов таких моделей обучают компактные головы под конкретную задачу, не трогая саму pLM; на этом принципе построены, помимо самого \textbf{DeepPeptide}~\cite{deeppeptide2023}, ещё \textbf{DeepNeuropePred}, \textbf{NeuroPred-PLM}, \textbf{SignalP~6}, а также модели оценки уже готовых пептидов: \textbf{pLM4ACE}~\cite{du2024plm4ace} определяет, ингибирует ли пептид АПФ, а \textbf{BPFun}~\cite{bpfun2025} и \textbf{DeepBP}~\cite{zhang2024deepbp} предсказывают его биоактивность. До-pLM \textbf{PeptideLocator}~\cite{mooney2013peptidelocator} -- ближайший по постановке предшественник самого DeepPeptide, но выдаёт лишь «тепловую карту» похожести остатка на часть биоактивного пептида, а не конкретные границы сегментов.

Из всех перечисленных моделей именно DeepPeptide решает задачу как посегментную разметку precursor-белка сразу с обеими метками (пептид/пропептид) и на момент публикации показывает лучшее качество, поэтому мы берём её отправной точкой, а не переизобретаем архитектуру заново. Наш вклад в этой части работы двоякий: переисправленная и переобученная реализация DeepPeptide на актуальных данных с задокументированной ошибкой метрики (раздел~\ref{sec:intro-flaws}), и систематическая проверка архитектурных взаимодействий (раздел~\ref{sec:results}) вместо однократного набора модификаций, отобранного и оценённого на одних и тех же данных.

\section{Постановка задачи}
\label{sec:problem}

Дан белок-предшественник в виде последовательности аминокислот
$$ x = (a_1, a_2, \dots, a_L), \qquad a_t \in \Omega, $$
где $\Omega$ -- алфавит из 20 аминокислот, $L$ -- длина последовательности. Нужно предсказать, какие непрерывные участки этой последовательности станут после расщепления зрелыми \textbf{пептидами}, а какие -- \textbf{пропептидами} (вспомогательными участками, которые при созревании белка отрезаются и удаляются и сами не являются конечным активным продуктом). Формально это сводится к присвоению каждой позиции метки
$$ y_t \in \{\textsf{None},\ \textsf{Peptide},\ \textsf{Propeptide}\}, $$
где непрерывные участки одинаковой метки образуют \textbf{сегменты}: типизированные пары (начало, конец).
$$
\underbrace{\texttt{MGK}}_{\textsf{None}}\;
\underbrace{\texttt{RALR}}_{\textsf{пропептид}}\;
\underbrace{\texttt{RSGK}}_{\textsf{пептид}}\;
\underbrace{\texttt{VR}}_{\textsf{None}}\;
\underbrace{\texttt{NL}}_{\textsf{пептид}}
$$

\subsection{Критерий совпадения: почему допуск, а не точное совпадение}

Предсказанные и истинные сегменты сопоставляются \emph{раздельно} для пептидов и пропептидов (пептид сравнивается только с пептидами и т.д.). Предсказанный сегмент засчитывается верным (true positive), если его начало и конец совпадают с началом и концом истинного сегмента того же типа с точностью до допуска $\pm\tau$ ($\tau=3$ остатка); непарный предсказанный сегмент считается ложным срабатыванием, непарный истинный -- пропуском. По этим счётчикам для каждого типа отдельно считаются точность, полнота и F1.

Почему выбран именно такой, посегментный критерий с допуском, а не строго позиционное совпадение (например, per-residue accuracy или IoU по маске)? Есть несколько причин.

Прежде всего, сама биологическая значимость локализации места разреза измеряется приемлемой близостью к истинному, а не точным попаданием в отдельную аминокислотную позицию: аннотация сайта расщепления в базах данных сама несёт экспериментальную или куraторскую неопределённость в несколько остатков, и предсказание, промахнувшееся на 1--2 остатка, для большинства практических применений (раздел~\ref{sec:intro-flaws}) неотличимо от точного попадания.

Кроме того, чисто позиционный критерий плохо сочетается с природой типичной ошибки модели в этой задаче. Единственная сдвинутая на несколько остатков граница по позиционной метрике размазывает штраф по всем затронутым позициям сегмента, но не даёт прямого ответа на практически более важный вопрос: нашла ли модель этот пептид вообще. Посегментный критерий с допуском отвечает на оба вопроса сразу -- найден ли сегмент (через P/R по числу совпавших сегментов) и насколько точно локализована его граница (через развёртку по $\tau$, раздел~\ref{sec:results}).

Наконец, тот же критерий, что и в оригинальной статье DeepPeptide ($\tau=3$, раздельно по типам), сохраняет прямую сопоставимость чисел с опубликованной работой, что было отдельной целью этого исследования (раздел~\ref{sec:eval-metric}).

\section{Оценка качества}
\label{sec:eval}

\subsection{Ошибка в эталонной метрике: как мы с ней поступили}
\label{sec:eval-metric}

В реализации функции \texttt{get\_counts\_for\_protein} (\texttt{src/utils/manuscript\_metrics.py}), которая считает описанный выше критерий $\pm\tau$, обнаружено переиспользование имени переменной цикла. Внешний цикл идёт по истинным сегментам, вложенный -- по предсказанным, и оба используют одно и то же имя индекса; при совпадении флаг «сматчено» ошибочно проставляется по индексу предсказанного сегмента, а не истинного. Эта ошибка не критична в среднем (частично компенсируется), но систематически занижает полноту и F1, особенно на белках, где предсказанных сегментов больше, чем истинных, и обнаружена также в оригинальном коде DeepPeptide.

Тихо исправлять метрику мы не стали: это сдвинуло бы абсолютные числа и нарушило сопоставимость с опубликованной работой. Вместо этого метрика ведётся в двух вариантах: опубликованный (ошибочный) сохранён для прямого сравнения, а исправленный используется как основной результат данной работы. Все числа в разделе~\ref{sec:results} -- исправленные; от исправления меняются только абсолютные значения, ранжирование моделей остаётся прежним.

\subsection{Протокол: от единичного разбиения к вложенной кросс-валидации}
\label{sec:eval-protocol}

Данные разбиваются на фолды с ограничением на межфолдовую гомологию инструментом GraphPart~\cite{graphpart2023} (раздел~\ref{sec:related}, приложение о датасете), дополнительно сбалансированные по кластерам мотивов, фланкирующих место разреза, что повторяет методику оригинальной статьи.

Первая версия протокола данной работы, на которой получены все эксперименты из приложения~\ref{sec:appendix}, разбивала данные на 7 фолдов с фиксированными ролями (train/dev/model-select/sealed). Это уже решало проблему (2) из раздела~\ref{sec:intro-flaws}, отделяя выбор архитектуры от итоговой оценки, но сама итоговая оценка оставалась точечной: одно случайное разбиение на роли, одна пара отложенных фолдов. На практике несколько проверенных на этом протоколе модификаций дали противоположный по знаку эффект на двух разных отложенных фолдах (например, эффект gated-адаптера был близок к нулю на одном фолде и составил $+0.046$ F1 на другом). Иначе говоря, единичного разбиения оказалось недостаточно, чтобы уверенно сказать, реален эффект или нет.

Поэтому в финальном протоколе данной работы используется \textbf{$5\times4$ вложенная кросс-валидация}. Берутся 5 внешних фолдов; для каждого внешнего фолда $o$ обучаются 4 модели, каждая -- на оставшихся трёх фолдах с одним из оставшихся четырёх фолдов $i\ne o$ в роли валидации (для выбора эпохи), и тестируется на $o$, который эта модель не видела ни на обучении, ни на валидации. Всего получается $5\times4=20$ обучений на модель. Оценка внешнего фолда -- среднее по его 4 внутренним моделям, итоговое качество -- среднее по 5 внешним фолдам, а разброс считается как стандартное отклонение именно по внешним фолдам, а не по всем 20 ячейкам сразу (4 внутренние модели одного внешнего фолда используют один и тот же тестовый набор и потому не являются независимыми повторами). Такой протокол дороже по вычислениям, чем единичное разбиение (раздел~\ref{sec:limitations}), зато каждый белок ровно один раз выступает тестовым, а итоговая оценка усредняется по всем пяти разбиениям, а не зависит от одного удачного или неудачного фолда.

\subsection{Неоднородность данных как источник шума}

Есть и отдельная причина, по которой единичное разбиение недостаточно надёжно: сама неоднородность данных. Белки в датасете относятся более чем к тысяче родов организмов, от человека и мыши до ядовитых беспозвоночных; пептидные и пропептидные сегменты сильно различаются по длине, а места разреза -- по фланкирующему мотиву. Фолды сбалансированы по гомологии и по мотиву разреза, но не по длине сегментов и не по составу организмов, поэтому на разных фолдах одни и те же модели могут заметно менять места в рейтинге, а отдельные семейства (например, яд пауков \textit{Cyriopagopus}) на конкретном фолде проседают на десятки пунктов F1. Усреднение по пяти внешним фолдам вложенной CV -- прямой ответ на эту неоднородность: эффект состава конкретного фолда усредняется, а не решает исход сравнения в одиночку.

\section{Предлагаемый метод}
\label{sec:method}

Обе предлагаемые модификации надстраиваются над одной и той же базовой схемой DeepPeptide (pLM $\to$ CNN--BiLSTM $\to$ CRF) и адресуют проблему (3) из раздела~\ref{sec:intro-flaws}, отсутствие позиционно-специфичных эмиссий, но с разным обоснованием: одна работает на уровне декодера (CRF), другая -- на уровне входного представления.

\subsection{Boundary-голова}

CRF-декодер DeepPeptide использует набор состояний, по которому сегмент любой допустимой длины проходит путь «начало $\to$ внутренние $\to$ конец»: первые несколько состояний сегмента всегда соответствуют его первым остаткам, а последние несколько -- его последним остаткам, независимо от длины всего сегмента. То есть сама структура состояний CRF уже различает «близко к границе» и «внутри», но в базовой модели эмиссии всех внутренних состояний одного типа идентичны, и эта позиционная информация никак не используется на уровне признаков.

Boundary-голова -- небольшой полносвязный блок (LayerNorm $\to$ линейный слой $\to$ GELU $\to$ линейный слой) поверх контекстных по-остаточных признаков (выход CNN--BiLSTM). Для каждой позиции и каждого типа сегмента она выдаёт три дополнительных числа: насколько эта позиция похожа на начало, на внутренность и на конец сегмента. Эти числа добавляются как поправки к эмиссиям соответствующих позиционно-фиксированных состояний CRF. На уровне идеи важна инициализация: последний линейный слой головы инициализируется нулями, поэтому в начале обучения модель в точности эквивалентна базовой (без головы) и вносит поправки только в той мере, в какой это снижает функцию потерь. Модификация встраивается в существующие CRF-эмиссии, а не заменяет их, и потому не может ухудшить оптимизацию с нуля.

\subsection{Адаптер}

Вторая модификация не трогает декодер, а работает на входе. Перед тем как контекстный признак (по-остаточный эмбеддинг pLM) попадает в CNN--BiLSTM, он проходит через небольшой обучаемый блок LayerNorm $\to$ Linear $\to$ GELU, который перепроецирует и, как правило, сжимает его размерность. Обоснование здесь другое, чем у boundary-головы: pLM-эмбеддинг обучен без учителя на общей задаче восстановления замаскированных остатков, а не на задаче локализации сайта протеолитического разреза, и распределение его признаков не подогнано ни под конкретную последующую задачу, ни под сравнительно небольшую по числу параметров CNN--BiLSTM-CRF голову. Лёгкий обучаемый адаптер перед основной сетью -- стандартный приём в transfer learning: он даёт модели возможность перенормировать и переупаковать признаки под задачу и под ёмкость последующей маленькой сети, не трогая параметры самой pLM.

По месту действия модификации независимы (одна меняет декодер, другая -- вход), поэтому проверяются и по отдельности, и совместно (раздел~\ref{sec:results}).

\section{Эксперименты и результаты}
\label{sec:results}

\Todo{Раздел ждёт финальных чисел по ESM-C 6B (раздел~\ref{sec:esmc-cv} ниже) и повторной вычитки под протокол $2\times2$: сетка модификаций (без / +boundary / +адаптер / +обе) на двух эмбеддингах (ESM-2, ESM-C 6B). Черновые числа по ESM-2 (полная сетка, $4\times20=80$ обучений):
база $0.556\pm0.028$ / +boundary $0.580\pm0.028$ ($+0.024$) / +адаптер $0.582\pm0.026$ ($+0.026$) / +обе $0.612\pm0.023$ ($+0.057$, почти аддитивно).
ESM-C 6B: не завершено на момент написания (база в процессе локально, +boundary/+адаптер/+обе поставлены в очередь на внешнем сервере).}

\subsection{Результаты на ESM-C 6B}
\label{sec:esmc-cv}
\Todo{Заполнить после завершения вычислений на внешнем сервере.}

\section{Обсуждение}
\label{sec:discussion}
\Todo{Три пункта по плану: (1) почему насыщенный эмбеддинг сам по себе не помогает без архитектурной надстройки; (2) почему трудно честно измерять результат на гомологичных сплитах (неоднородность данных, раздел~\ref{sec:eval-protocol}); (3) почему адаптер вообще может помогать (гипотеза про перенормировку pLM-признаков под небольшую последующую сеть).}

\section{Ограничения}
\label{sec:limitations}
\Todo{Пункты: (1) нет отдельной запечатанной выборки сверх структуры вложенной CV и почему это трудно решить полностью; (2) часть идей (Ankh, LoRA-дообучение, контрастное обучение) не проверена на честном протоколе; (3) высокая вычислительная стоимость честного протокола -- одна ячейка CV $\approx$7--11 часов, полная сетка $\approx$150--200 GPU-часов.}

\section{Заключение}
\label{sec:conclusion}
\Todo{Сжатый пересказ обсуждения: (1) интеракция ESM-C 6B $\times$ boundary-голова; (2) адаптер как самостоятельный подтверждённый ингредиент; (3) чистота протокола как отдельная ценность работы.}

\appendix
\section{Приложение: более ранние наблюдения и материалы для воспроизводимости}
\label{sec:appendix}
\Todo{Перенести сюда полный перечень экспериментов, не проверенных на честном протоколе (3Di, bond-лосс, телескопический CRF, словарь Aho--Corasick, структурные признаки AlphaFold/AFToolkit, Ankh, LoRA), явно пометив их как наблюдения, не доказательства. Плюс: таблица гиперпараметров по конфигурациям, полные числа по всем 20 ячейкам каждой модели (не только среднее по внешним фолдам), раздел о воспроизводимости (код, датасет, разбиение).}

\newpage
\bibliographystyle{plain}
\bibliography{refs}

\end{document}
